\begin{theorem}[折叠规范群的 Stirling--Bernoulli 层级：由折叠数据内生读出全部 Bernoulli 系数]\label{thm:fold-bin-gauge-constant-stirling-bernoulli-hierarchy}
沿用推论 \ref{cor:fold-bin-recover-2pi} 的 \(C_m\) 定义，并记 \(\log C_m\) 为其自然对数。
则存在完整渐近展开：对任意固定整数 \(R\ge 1\)，当 \(m\to\infty\) 时
\[
\boxed{\;
\log C_m
=\log(2\pi)+
\sum_{r=1}^{R}
\frac{B_{2r}}{r(2r-1)}
\Bigl(\varphi^{-1}+\varphi^{2r-3}\Bigr)
\Bigl(\frac{\varphi}{2}\Bigr)^{(2r-1)m}
\;+\;
O\!\Bigl(\Bigl(\frac{\varphi}{2}\Bigr)^{(2R+1)m}\Bigr)\; } ,
\]
其中 \(B_{2r}\) 为 Bernoulli 数。
因此对每个 \(r\ge 1\)，Bernoulli 系数 \(B_{2r}\) 可由折叠数据序列 \(\{\log C_m\}_{m\ge 1}\) 通过逐阶消元在极限中唯一恢复；
并且结合标准恒等式
\[
B_{2r}=(-1)^{r-1}\frac{2(2r)!}{(2\pi)^{2r}}\zeta(2r),
\]
可将全部 \(\zeta(2r)\) 视作该折叠规范常数的渐近谱信息。
\end{theorem}
\begin{proof}
采用 Stirling--Bernoulli 渐近展开（作为外部工具）：对任意固定 \(R\ge 1\) 与整数 \(d\ge 1\) 有
\[
\log(d!)=d\log d-d+\frac12\log(2\pi d)+
\sum_{r=1}^{R}\frac{B_{2r}}{2r(2r-1)}\,d^{-(2r-1)}
\;+\;
O\!\bigl(d^{-(2R+1)}\bigr).
\]
将其应用到 \(d=d^{\mathrm{bin}}_m(w)\) 并对 \(w\in X_m\) 求和，再除以 \(2^m\)，并按
\(
g_m:=2^{-m}\sum_w\log(d^{\mathrm{bin}}_m(w)!)
\)
与
\(
\kappa_m:=2^{-m}\sum_w d^{\mathrm{bin}}_m(w)\log d^{\mathrm{bin}}_m(w)
\)
整理，可得
\[
g_m
=\kappa_m-1+\frac{1}{2^{m+1}}\sum_{w\in X_m}\log(2\pi d^{\mathrm{bin}}_m(w))
\;+\;
\sum_{r=1}^{R}\frac{B_{2r}}{2r(2r-1)}\cdot
\frac{1}{2^{m}}\sum_{w\in X_m} d^{\mathrm{bin}}_m(w)^{-(2r-1)}
\;+\;
O\!\left(\frac{1}{2^{m}}\sum_{w\in X_m} d^{\mathrm{bin}}_m(w)^{-(2R+1)}\right).
\]
将
\(
\sum_w\log(2\pi d)=|X_m|\log(2\pi)+\sum_w\log d
\)
代回，并代入 \(\log C_m\) 的定义（推论 \ref{cor:fold-bin-recover-2pi}），主项严格抵消，仅留下 Bernoulli 级数的均匀平均：
\[
\log C_m
=\log(2\pi)+
\sum_{r=1}^{R}\frac{B_{2r}}{r(2r-1)}\cdot
\frac{1}{|X_m|}\sum_{w\in X_m} d^{\mathrm{bin}}_m(w)^{-(2r-1)}
\;+\;
O\!\left(\frac{1}{|X_m|}\sum_{w\in X_m} d^{\mathrm{bin}}_m(w)^{-(2R+1)}\right).
\]
\par
由定理 \ref{thm:fold-bin-two-state-asymptotic} 的一致渐近
\(
d^{\mathrm{bin}}_m(w)=\varphi^{-w_m}(2/\varphi)^m+O(1)
\)
以及末位分块 \(\#\{w\in X_m:w_m=0\}=F_{m+1}\)、\(\#\{w\in X_m:w_m=1\}=F_m\) 与 \(F_{m+1}/F_{m+2}\to\varphi^{-1}\)、\(F_m/F_{m+2}\to\varphi^{-2}\)，可得对每个固定 \(r\ge 1\) 有
\[
\frac{1}{|X_m|}\sum_{w\in X_m} d^{\mathrm{bin}}_m(w)^{-(2r-1)}
=\Bigl(\varphi^{-1}+\varphi^{2r-3}\Bigr)\Bigl(\frac{\varphi}{2}\Bigr)^{(2r-1)m}
\;+\;
O\!\Bigl(\Bigl(\frac{\varphi}{2}\Bigr)^{(2r+1)m}\Bigr).
\]
同理，
\(
|X_m|^{-1}\sum_w d^{-(2R+1)}
=O((\varphi/2)^{(2R+1)m})
\)。
代回即得所示展开与余项阶。最后两句关于“逐阶消元恢复 \(B_{2r}\)”与 \(\zeta(2r)\) 的编码为标准推论：由于各阶指数 \((\varphi/2)^{(2r-1)m}\) 相互分离，可通过依次去除低阶项并令 \(m\to\infty\) 唯一读出每个系数。证毕。
\end{proof}

\begin{proposition}[Bernoulli 系数的显式抽取算子：指数分离给出无外参的极限读出]\label{prop:fold-bin-gauge-bernoulli-extraction-operator}
\leanverified{paper\_fold\_bin\_gauge\_bernoulli\_extraction\_operator}
沿用定理 \ref{thm:fold-bin-gauge-constant-stirling-bernoulli-hierarchy} 的记号，并令
\[
q:=\frac{\varphi}{2}\in(0,1),\qquad
L_m:=\log C_m-\lim_{k\to\infty}\log C_k=\log C_m-\log(2\pi).
\]
递归定义系数读出：
\[
a_1:=\lim_{m\to\infty}q^{-m}L_m,
\]
以及对每个 \(r\ge 2\),
\[
a_r:=\lim_{m\to\infty} q^{-(2r-1)m}\left(L_m-\sum_{j=1}^{r-1} a_j\,q^{(2j-1)m}\right).
\]
则每个极限 \(a_r\) 均存在，并满足
\[
a_r=\frac{B_{2r}}{r(2r-1)}\Bigl(\varphi^{-1}+\varphi^{2r-3}\Bigr).
\]
从而全部 \(B_{2r}\)（并进而全部 \(\zeta(2r)\)）可由折叠输出序列 \(\{\log C_m\}_{m\ge 1}\) 逐阶确定性恢复；该恢复仅使用 \(C_m\) 序列本身及其极限 \(\lim_{m\to\infty}\log C_m\)，不依赖外部拟合参数。
\end{proposition}
\begin{proof}
由定理 \ref{thm:fold-bin-gauge-constant-stirling-bernoulli-hierarchy} 的展开，对任意固定 \(R\ge 1\) 有
\[
L_m=\sum_{r=1}^{R}\alpha_r\,q^{(2r-1)m}+O\!\bigl(q^{(2R+1)m}\bigr),
\qquad m\to\infty,
\]
其中
\[
\alpha_r:=\frac{B_{2r}}{r(2r-1)}\Bigl(\varphi^{-1}+\varphi^{2r-3}\Bigr).
\]
取 \(R=1\) 得
\(
q^{-m}L_m=\alpha_1+O(q^{2m})
\)，故极限 \(a_1=\lim_{m\to\infty}q^{-m}L_m\) 存在且 \(a_1=\alpha_1\)。
\par
设对某个 \(r\ge 2\) 已有 \(a_j=\alpha_j\)（\(1\le j\le r-1\)）。对定理展开取 \(R=r\)，移项并代入归纳假设，得
\[
L_m-\sum_{j=1}^{r-1}a_j\,q^{(2j-1)m}
=\alpha_r\,q^{(2r-1)m}+O\!\bigl(q^{(2r+1)m}\bigr).
\]
两边乘以 \(q^{-(2r-1)m}\) 得
\[
q^{-(2r-1)m}\left(L_m-\sum_{j=1}^{r-1}a_j\,q^{(2j-1)m}\right)
=\alpha_r+O(q^{2m}),
\]
从而极限 \(a_r\) 存在且 \(a_r=\alpha_r\)。最后由
\(
B_{2r}=(-1)^{r-1}\frac{2(2r)!}{(2\pi)^{2r}}\zeta(2r)
\)
即可逐阶恢复全部偶值 \(\zeta(2r)\)。证毕。
\end{proof}

\endinput
